% Options for packages loaded elsewhere
\PassOptionsToPackage{unicode}{hyperref}
\PassOptionsToPackage{hyphens}{url}
\documentclass[
  11pt,
]{article}
\usepackage{xcolor}
\usepackage[margin=1in]{geometry}
\usepackage{amsmath,amssymb}
\setcounter{secnumdepth}{5}
\usepackage{iftex}
\ifPDFTeX
  \usepackage[T1]{fontenc}
  \usepackage[utf8]{inputenc}
  \usepackage{textcomp} % provide euro and other symbols
\else % if luatex or xetex
  \usepackage{unicode-math} % this also loads fontspec
  \defaultfontfeatures{Scale=MatchLowercase}
  \defaultfontfeatures[\rmfamily]{Ligatures=TeX,Scale=1}
\fi
\usepackage{lmodern}
\ifPDFTeX\else
  % xetex/luatex font selection
\fi
% Use upquote if available, for straight quotes in verbatim environments
\IfFileExists{upquote.sty}{\usepackage{upquote}}{}
\IfFileExists{microtype.sty}{% use microtype if available
  \usepackage[]{microtype}
  \UseMicrotypeSet[protrusion]{basicmath} % disable protrusion for tt fonts
}{}
\makeatletter
\@ifundefined{KOMAClassName}{% if non-KOMA class
  \IfFileExists{parskip.sty}{%
    \usepackage{parskip}
  }{% else
    \setlength{\parindent}{0pt}
    \setlength{\parskip}{6pt plus 2pt minus 1pt}}
}{% if KOMA class
  \KOMAoptions{parskip=half}}
\makeatother
\usepackage{color}
\usepackage{fancyvrb}
\newcommand{\VerbBar}{|}
\newcommand{\VERB}{\Verb[commandchars=\\\{\}]}
\DefineVerbatimEnvironment{Highlighting}{Verbatim}{commandchars=\\\{\}}
% Add ',fontsize=\small' for more characters per line
\usepackage{framed}
\definecolor{shadecolor}{RGB}{248,248,248}
\newenvironment{Shaded}{\begin{snugshade}}{\end{snugshade}}
\newcommand{\AlertTok}[1]{\textcolor[rgb]{0.94,0.16,0.16}{#1}}
\newcommand{\AnnotationTok}[1]{\textcolor[rgb]{0.56,0.35,0.01}{\textbf{\textit{#1}}}}
\newcommand{\AttributeTok}[1]{\textcolor[rgb]{0.13,0.29,0.53}{#1}}
\newcommand{\BaseNTok}[1]{\textcolor[rgb]{0.00,0.00,0.81}{#1}}
\newcommand{\BuiltInTok}[1]{#1}
\newcommand{\CharTok}[1]{\textcolor[rgb]{0.31,0.60,0.02}{#1}}
\newcommand{\CommentTok}[1]{\textcolor[rgb]{0.56,0.35,0.01}{\textit{#1}}}
\newcommand{\CommentVarTok}[1]{\textcolor[rgb]{0.56,0.35,0.01}{\textbf{\textit{#1}}}}
\newcommand{\ConstantTok}[1]{\textcolor[rgb]{0.56,0.35,0.01}{#1}}
\newcommand{\ControlFlowTok}[1]{\textcolor[rgb]{0.13,0.29,0.53}{\textbf{#1}}}
\newcommand{\DataTypeTok}[1]{\textcolor[rgb]{0.13,0.29,0.53}{#1}}
\newcommand{\DecValTok}[1]{\textcolor[rgb]{0.00,0.00,0.81}{#1}}
\newcommand{\DocumentationTok}[1]{\textcolor[rgb]{0.56,0.35,0.01}{\textbf{\textit{#1}}}}
\newcommand{\ErrorTok}[1]{\textcolor[rgb]{0.64,0.00,0.00}{\textbf{#1}}}
\newcommand{\ExtensionTok}[1]{#1}
\newcommand{\FloatTok}[1]{\textcolor[rgb]{0.00,0.00,0.81}{#1}}
\newcommand{\FunctionTok}[1]{\textcolor[rgb]{0.13,0.29,0.53}{\textbf{#1}}}
\newcommand{\ImportTok}[1]{#1}
\newcommand{\InformationTok}[1]{\textcolor[rgb]{0.56,0.35,0.01}{\textbf{\textit{#1}}}}
\newcommand{\KeywordTok}[1]{\textcolor[rgb]{0.13,0.29,0.53}{\textbf{#1}}}
\newcommand{\NormalTok}[1]{#1}
\newcommand{\OperatorTok}[1]{\textcolor[rgb]{0.81,0.36,0.00}{\textbf{#1}}}
\newcommand{\OtherTok}[1]{\textcolor[rgb]{0.56,0.35,0.01}{#1}}
\newcommand{\PreprocessorTok}[1]{\textcolor[rgb]{0.56,0.35,0.01}{\textit{#1}}}
\newcommand{\RegionMarkerTok}[1]{#1}
\newcommand{\SpecialCharTok}[1]{\textcolor[rgb]{0.81,0.36,0.00}{\textbf{#1}}}
\newcommand{\SpecialStringTok}[1]{\textcolor[rgb]{0.31,0.60,0.02}{#1}}
\newcommand{\StringTok}[1]{\textcolor[rgb]{0.31,0.60,0.02}{#1}}
\newcommand{\VariableTok}[1]{\textcolor[rgb]{0.00,0.00,0.00}{#1}}
\newcommand{\VerbatimStringTok}[1]{\textcolor[rgb]{0.31,0.60,0.02}{#1}}
\newcommand{\WarningTok}[1]{\textcolor[rgb]{0.56,0.35,0.01}{\textbf{\textit{#1}}}}
\usepackage{longtable,booktabs,array}
\usepackage{calc} % for calculating minipage widths
% Correct order of tables after \paragraph or \subparagraph
\usepackage{etoolbox}
\makeatletter
\patchcmd\longtable{\par}{\if@noskipsec\mbox{}\fi\par}{}{}
\makeatother
% Allow footnotes in longtable head/foot
\IfFileExists{footnotehyper.sty}{\usepackage{footnotehyper}}{\usepackage{footnote}}
\makesavenoteenv{longtable}
\usepackage{graphicx}
\makeatletter
\newsavebox\pandoc@box
\newcommand*\pandocbounded[1]{% scales image to fit in text height/width
  \sbox\pandoc@box{#1}%
  \Gscale@div\@tempa{\textheight}{\dimexpr\ht\pandoc@box+\dp\pandoc@box\relax}%
  \Gscale@div\@tempb{\linewidth}{\wd\pandoc@box}%
  \ifdim\@tempb\p@<\@tempa\p@\let\@tempa\@tempb\fi% select the smaller of both
  \ifdim\@tempa\p@<\p@\scalebox{\@tempa}{\usebox\pandoc@box}%
  \else\usebox{\pandoc@box}%
  \fi%
}
% Set default figure placement to htbp
\def\fps@figure{htbp}
\makeatother
\setlength{\emergencystretch}{3em} % prevent overfull lines
\providecommand{\tightlist}{%
  \setlength{\itemsep}{0pt}\setlength{\parskip}{0pt}}
\let\maketitle\relax
\usepackage{float}
\usepackage{bookmark}
\IfFileExists{xurl.sty}{\usepackage{xurl}}{} % add URL line breaks if available
\urlstyle{same}
\hypersetup{
  hidelinks,
  pdfcreator={LaTeX via pandoc}}

\author{}
\date{\vspace{-2.5em}}

\begin{document}

\thispagestyle{empty}
\begin{center}
\vspace*{0.5cm}
{\Large \textbf{Newcastle University Business School}}\\[6pt]
{\large ISO8005 — Web \& Social Media Analytics}\\[18pt]

\rule{\linewidth}{0.4pt}\\[6pt]
{\LARGE \textbf{Assignment Cover Sheet}}\\[6pt]
\rule{\linewidth}{0.4pt}\\[18pt]
\end{center}

\begin{tabular}{p{4.5cm} p{10cm}}
\textbf{Module Code:}     & ISO8005 \\
\textbf{Module Name:}     & Web \& Social Media Analytics \\
\textbf{Assignment:}      & Tyne Artisan Bakery — From Oven to Online \\
\textbf{Academic Year:}   & 2025–26 \\
\textbf{Student Number:}  & [insert student number] \\
\textbf{Submission Date:} & May 2026 \\
\end{tabular}

\vspace{0.8cm}

\textbf{Assignment Declarations} \vspace{0.2cm}

\textit{Declaration of Original Work} --- I certify that this submission
is entirely my own work and has not been copied, reproduced, or
completed with unauthorised assistance. All sources used have been
appropriately acknowledged. \hfill \(\boxtimes\)

\vspace{0.4cm}

\textit{Declaration of AI Understanding} --- I certify that I have read
and understood the AI for Students Academic Integrity Checklist.
\hfill \(\boxtimes\)

\vspace{0.4cm}

\textit{Declaration of Word Count} --- Total word count for this
assignment is: {[}insert word count{]} words (excluding cover page,
executive summary, contents, references, appendices and visuals, per the
brief).

\vspace{0.4cm}

\textit{Declaration of Development Work Retention} --- All development
works, scripts, RData snapshots and figures are retained in the project
repository for a minimum of three months after submission and are
available for scrutiny within five working days of request.
\hfill \(\boxtimes\)

\vspace{1.2cm}
\rule{\linewidth}{0.4pt}

\textbackslash{[}4pt{]}
\textbf{Project Access Points}\textbackslash{[}6pt{]}

\rule{\linewidth}{0.4pt}

\begin{tabular}{p{4.5cm} p{10cm}}
\textbf{Interactive dashboard:} & \texttt{output/TAB\_Dashboard.html} (offline, self-contained HTML) \\
\textbf{Source repository:}     & \texttt{https://github.com/<your-handle>/iso8005-tab-report} \emph{(placeholder)} \\
\end{tabular}

\vspace{0.4cm}

\small \emph{The interactive dashboard renders 11 KPI panels from the same RData snapshots used in this report (Figure 16). The GitHub repository contains the Rmd, R scripts, dataset, figures and the offline dashboard.}

\newpage

\thispagestyle{empty}
\begin{center}
\vspace*{2cm}
{\LARGE \textbf{From Oven to Online}}\\[10pt]
{\Large A Data-Led, AI-Augmented Go-to-Market Plan for}\\[4pt]
{\Large \textbf{Tyne Artisan Bakery}}\\[18pt]
{\large \emph{Anchored on Monthly Active Local Customers (MALC) as the North Star metric}}\\[24pt]
\rule{0.6\linewidth}{0.4pt}\\[14pt]
{\large Student Number: \texttt{[insert student number]}}\\[6pt]
{\large May 2026}\\[6pt]
{\normalsize ISO8005 --- Web \& Social Media Analytics}\\
{\normalsize Newcastle University Business School}
\end{center}
\vfill

\newpage

\tableofcontents

\newpage

\section{Executive Summary}\label{executive-summary}

\begin{quote}
\emph{``Growth is not only about acquiring more customers; it is about
repeatedly creating measurable value for the right customers.''}
\end{quote}

Tyne Artisan Bakery (TAB) is a Newcastle artisan bakery with a strong
in-store reputation but no digital presence. This report sets out a
12--18 month, data-led and AI-augmented go-to-market plan anchored on a
single North Star metric --- \textbf{Monthly Active Local Customers
(MALC)} --- and grounded in TAB's two-year dataset (5,942 orders, 4,179
unique customers, £120,835 revenue, Mar 2024 -- Feb 2026).

\begin{longtable}[]{@{}
  >{\raggedright\arraybackslash}p{(\linewidth - 4\tabcolsep) * \real{0.2148}}
  >{\raggedright\arraybackslash}p{(\linewidth - 4\tabcolsep) * \real{0.2815}}
  >{\raggedright\arraybackslash}p{(\linewidth - 4\tabcolsep) * \real{0.5037}}@{}}
\caption{Executive summary of TAB's commercial baseline}\tabularnewline
\toprule\noalign{}
\begin{minipage}[b]{\linewidth}\raggedright
Metric
\end{minipage} & \begin{minipage}[b]{\linewidth}\raggedright
Value
\end{minipage} & \begin{minipage}[b]{\linewidth}\raggedright
Strategic meaning
\end{minipage} \\
\midrule\noalign{}
\endfirsthead
\toprule\noalign{}
\begin{minipage}[b]{\linewidth}\raggedright
Metric
\end{minipage} & \begin{minipage}[b]{\linewidth}\raggedright
Value
\end{minipage} & \begin{minipage}[b]{\linewidth}\raggedright
Strategic meaning
\end{minipage} \\
\midrule\noalign{}
\endhead
\bottomrule\noalign{}
\endlastfoot
Orders analysed & 5,942 & Sufficient order base for customer, channel
and conversion analysis \\
Unique customers & 4,179 & Shows a sizeable customer pool for
retention-led growth \\
Total revenue & £120,835 & Confirms commercial potential, but not yet a
mature digital engine \\
Average order value & £20.34 & Requires careful acquisition-cost control
before scaling paid media \\
North Star metric & Monthly Active Local Customers (MALC) & Connects
digital growth with repeat local customer engagement \\
Recommended planning horizon & 12--18 months & Allows phased MVP launch,
testing, optimisation and scale-up \\
\end{longtable}

\textbf{Four headline findings drive the plan:}

\begin{enumerate}
\def\labelenumi{\arabic{enumi}.}
\tightlist
\item
  \textbf{Paid media is unprofitable.} Only Google Search clears the
  1.0\(\times\) ROAS break-even (1.14\(\times\)); Meta 0.82\(\times\),
  TikTok 0.90\(\times\), Other 0.62\(\times\). CAC (£36--£52) exceeds
  AOV (£20.34) --- TAB must lead with owned channels (email,
  subscriptions, local SEO).
\item
  \textbf{North East is the launch market.} It already delivers 20.1\%
  of revenue, the natural catchment for local delivery and
  click-and-collect.
\item
  \textbf{The hero-banner A/B test is decisive.} Variant B converts at
  3.65\% vs 2.86\% --- a +27.5\% uplift, z = 6.90, p \textless{} 0.001
  --- and should ship from day one.
\item
  \textbf{Mobile-first, subscription-led.} 69\% of orders are mobile and
  average monthly churn is 4.5\%; an AI recommendation layer plus a
  subscription engine compounds retention.
\end{enumerate}

\textbf{Recommended solution:} a Shopify Basic MVP with a thin AI layer
(LLM recommender, Prophet demand forecasting, GenAI content studio for
menu copy and alt text). At a 12-month spend of \textbf{£11,500}, the
Base scenario (+4\%/mo MALC) lifts MALC from \textbf{430 (Feb 2026) to
\(\approx\) 688 (Feb 2027)}; High adds a further \(\approx\) 30\%. A
live leadership dashboard summarising all KPIs accompanies this report
and is listed on the cover sheet.

\newpage

\section{Introduction}\label{introduction}

\subsection{Business Context}\label{business-context}

TAB is a microenterprise: founder-led, fewer than ten staff, bounded by
the physical reach of an in-store counter and weekend market. The owners
want to launch e-commerce (local delivery and click-and-collect) and
build a coherent social presence from scratch, with modest budgets and
limited technical capacity.

The deeper strategic problem is the \textbf{absence of a measurable
digital customer-acquisition and -retention system}. Without
instrumented funnels and a single agreed metric, every marketing
investment becomes a bet rather than an experiment. This report
sequences (i) the metric, (ii) the MVP that generates the data, and
(iii) the analytics and forecasting that turn the data into
reallocation.

\subsection{Methodology}\label{methodology}

The analysis follows a CRISP-DM workflow (Wirth \& Hipp, 2000) with
reproducible-research practice (Wickham \& Bryan, 2023). Four stages run
end-to-end via \texttt{source("scripts/RUN\_ALL.R")}:

\textbf{Stage 1 --- Data Understanding} (Script 01). Load the 7-sheet
workbook via \texttt{readxl}; harmonise types with
\texttt{janitor::clean\_names()} and \texttt{lubridate}; profile null
counts and date ranges.

\textbf{Stage 2 --- Metric Modelling} (Script 02). Compute 15+ KPIs
grouped by month/channel/region/device/variant. Key formulas:

\begin{longtable}[]{@{}
  >{\raggedright\arraybackslash}p{(\linewidth - 4\tabcolsep) * \real{0.3333}}
  >{\raggedright\arraybackslash}p{(\linewidth - 4\tabcolsep) * \real{0.3333}}
  >{\raggedright\arraybackslash}p{(\linewidth - 4\tabcolsep) * \real{0.3333}}@{}}
\toprule\noalign{}
\begin{minipage}[b]{\linewidth}\raggedright
KPI
\end{minipage} & \begin{minipage}[b]{\linewidth}\raggedright
Formula
\end{minipage} & \begin{minipage}[b]{\linewidth}\raggedright
Source
\end{minipage} \\
\midrule\noalign{}
\endhead
\bottomrule\noalign{}
\endlastfoot
CVR & \texttt{orders\ /\ sessions} & \texttt{monthly\_marketing} \\
AOV & \texttt{revenue\ /\ orders} & \texttt{orders} \\
ROAS & \texttt{attributed\_revenue\ /\ spend} &
\texttt{paid\_campaigns} \\
CAC & \texttt{spend\ /\ new\_customers} & \texttt{paid\_campaigns} \\
MALC & \texttt{n\_distinct(customer\_id)} per month & \texttt{orders} \\
Churn & \texttt{subscription\_cancels\ /\ active\_subscriptions} &
\texttt{monthly\_marketing} \\
CLV & \texttt{(AOV\ *\ orders\_per\_cust\ *\ margin)\ /\ churn} &
derived \\
\end{longtable}

\textbf{Stage 3 --- Statistical Inference}. Two-proportion z-test for
the A/B uplift (Agresti, 2018); compound-growth scenario forecasting;
simplified Pareto/NBD CLV (Fader, Hardie \& Lee, 2005).

\textbf{Stage 4 --- Decision Design}. Every insight is mapped to a
specific action with an expected MALC, ROAS or CAC impact.

\newpage

\section{TAB Web and Social Media Analysis \&
Plan}\label{tab-web-and-social-media-analysis-plan}

\subsection{Problem Framing, North Star and SMART
Objectives}\label{problem-framing-north-star-and-smart-objectives}

\textbf{MALC definition:} \emph{the count of unique customers (by
\texttt{customer\_id}) who place \(\geq\) 1 online order,
click-and-collect order, or active bread subscription within a calendar
month, weighted to the North East catchment.} MALC satisfies the four
good-North-Star tests in Croll \& Yoskovitz (2013): aligned to value,
leading, actionable, understandable.

\begin{figure}[H]

{\centering \includegraphics[width=0.8\linewidth]{figures/fig1_growth_logic} 

}

\caption{MALC growth logic — funnel-to-North-Star mapping.}\label{fig:unnamed-chunk-1}
\end{figure}

\textbf{SMART objectives (12-month):}

\begin{longtable}[]{@{}
  >{\raggedright\arraybackslash}p{(\linewidth - 6\tabcolsep) * \real{0.2500}}
  >{\raggedright\arraybackslash}p{(\linewidth - 6\tabcolsep) * \real{0.2500}}
  >{\raggedright\arraybackslash}p{(\linewidth - 6\tabcolsep) * \real{0.2500}}
  >{\raggedright\arraybackslash}p{(\linewidth - 6\tabcolsep) * \real{0.2500}}@{}}
\toprule\noalign{}
\begin{minipage}[b]{\linewidth}\raggedright
Objective
\end{minipage} & \begin{minipage}[b]{\linewidth}\raggedright
KPI
\end{minipage} & \begin{minipage}[b]{\linewidth}\raggedright
Target
\end{minipage} & \begin{minipage}[b]{\linewidth}\raggedright
Deadline
\end{minipage} \\
\midrule\noalign{}
\endhead
\bottomrule\noalign{}
\endlastfoot
Grow local customer base & MALC & +60\% (430 \(\rightarrow\)
\$\approx\$690) & Feb 2027 \\
Improve site conversion & CVR & \(\geq\) 3.0\% & End Q3 2026 \\
Build repeat revenue & Active subscriptions & \(\geq\) 700 & End Q4
2026 \\
Optimise paid efficiency & Blended ROAS & \(\geq\) 1.5\(\times\) & End
Q2 2026 \\
Build owned channel & Email list size & \(\geq\) 2,500 sign-ups & End Q2
2026 \\
\end{longtable}

\subsection{E-commerce MVP \& UX --- AI-Augmented
Shopify}\label{e-commerce-mvp-ux-ai-augmented-shopify}

\textbf{Platform.} Shopify Basic delivers PCI-compliant payments, GDPR
cookie tooling, a subscription app ecosystem (Recharge), email
(Klaviyo), and a mobile-first PWA theme out of the box (Shopify, 2025)
--- within the founder's budget and technical capacity. WooCommerce
shifts security/infra onto the owner; custom build is out of scope.

\textbf{AI layer (Figure 2).} LLM recommender, Prophet demand
forecasting (Taylor \& Letham, 2018), GBM churn-risk scoring feeding
Klaviyo win-back, and a GenAI content studio for menu copy, alt text and
captions --- all under founder review.

\begin{figure}[H]

{\centering \includegraphics[width=0.85\linewidth]{figures/fig8_ai_mvp_architecture} 

}

\caption{AI-augmented MVP architecture.}\label{fig:unnamed-chunk-2}
\end{figure}

\textbf{Information architecture (Figure 3).} Three-click depth; trust
artefacts (privacy, returns, allergens, accessibility) clustered in the
footer --- pattern shown to cut checkout abandonment (Baymard, 2024).

\begin{figure}[H]

{\centering \includegraphics[width=0.85\linewidth]{figures/fig10_site_map} 

}

\caption{TAB information architecture (site map).}\label{fig:unnamed-chunk-3}
\end{figure}

\textbf{UX wireframes (Figure 4).} Mobile Home / PDP / one-page
Checkout. Apple Pay / Google Pay above the fold cut purchase paths by
35--45\% in DTC food (Stripe, 2025).

\begin{figure}[H]

{\centering \includegraphics[width=0.85\linewidth]{figures/fig9_ux_wireframes} 

}

\caption{UX wireframes: mobile-first MVP home, product and checkout screens.}\label{fig:unnamed-chunk-4}
\end{figure}

\subsection{Channel \& Content Strategy}\label{channel-content-strategy}

\textbf{Personas (Figure 5).} P1 Weekly Bread Buyer (loyal subscriber);
P2 Celebration Cake Buyer (event-driven, high AOV); P3 Weekend Treat
Buyer (impulse mobile shopper).

\begin{figure}[H]

{\centering \includegraphics[width=0.85\linewidth]{figures/fig7_persona_funnel} 

}

\caption{Persona x Funnel content matrix.}\label{fig:unnamed-chunk-5}
\end{figure}

\textbf{Budget split (Q1):} Google 60 / Meta 25 / TikTok 15 --- inverted
from the typical DTC-food default because Google is the only channel
currently above 1.0\(\times\) ROAS (§3.4). Revisited quarterly using the
rule in §3.6.

\textbf{Content pillars.} Craft \& process (Reels/TikTok
3\(\times\)/wk), Weekly menu drop (email + Story), Allergen \& education
(monthly blog), Local community (monthly micro-influencer). KPIs in each
row link to a SMART objective.

\textbf{Experimentation.} Continuous testing programme to MDE 10\%,
\(\alpha\) = 0.05, power 0.80 (Kohavi, Tang \& Xu, 2020). Hero-banner
test evidence is shown in Figure 12; it is the first experiment to ship,
and the next queue includes (a) free-delivery threshold £25 vs £35, (b)
subscription discount 5\% vs 10\%, (c) PDP image-first vs reviews-first.

\subsection{Data Analysis --- Key
Insights}\label{data-analysis-key-insights}

\subsubsection{MALC and Revenue Trend}\label{malc-and-revenue-trend}

The blended MALC and revenue series shows a noisy but clearly positive
trajectory: from 92 customers (Mar 2024) to 430 (Feb 2026).
Trailing-12-month compound monthly growth is \(\approx\) 4.0\%, the
anchor for the Base forecast.

\begin{figure}[H]

{\centering \includegraphics[width=0.85\linewidth]{figures/fig2_malc_revenue_trend} 

}

\caption{MALC and revenue trend.}\label{fig:unnamed-chunk-6}
\end{figure}

\subsubsection{Paid Channel Efficiency}\label{paid-channel-efficiency}

Only Google clears break-even (Figure 7); blended CAC (£36--£52) exceeds
AOV on every paid platform. TAB cannot grow profitably through paid
alone --- lead with owned channels.

\begin{figure}[H]

{\centering \includegraphics[width=0.85\linewidth]{figures/fig3_channel_roas_cac} 

}

\caption{Paid channel efficiency (ROAS vs CAC).}\label{fig:unnamed-chunk-7}
\end{figure}

\begin{longtable}[]{@{}
  >{\raggedright\arraybackslash}p{(\linewidth - 14\tabcolsep) * \real{0.1250}}
  >{\raggedright\arraybackslash}p{(\linewidth - 14\tabcolsep) * \real{0.1250}}
  >{\raggedright\arraybackslash}p{(\linewidth - 14\tabcolsep) * \real{0.1250}}
  >{\raggedright\arraybackslash}p{(\linewidth - 14\tabcolsep) * \real{0.1250}}
  >{\raggedright\arraybackslash}p{(\linewidth - 14\tabcolsep) * \real{0.1250}}
  >{\raggedright\arraybackslash}p{(\linewidth - 14\tabcolsep) * \real{0.1250}}
  >{\raggedright\arraybackslash}p{(\linewidth - 14\tabcolsep) * \real{0.1250}}
  >{\raggedright\arraybackslash}p{(\linewidth - 14\tabcolsep) * \real{0.1250}}@{}}
\toprule\noalign{}
\begin{minipage}[b]{\linewidth}\raggedright
Channel
\end{minipage} & \begin{minipage}[b]{\linewidth}\raggedright
Spend
\end{minipage} & \begin{minipage}[b]{\linewidth}\raggedright
Clicks
\end{minipage} & \begin{minipage}[b]{\linewidth}\raggedright
CTR
\end{minipage} & \begin{minipage}[b]{\linewidth}\raggedright
CPC
\end{minipage} & \begin{minipage}[b]{\linewidth}\raggedright
Conv.
\end{minipage} & \begin{minipage}[b]{\linewidth}\raggedright
ROAS
\end{minipage} & \begin{minipage}[b]{\linewidth}\raggedright
CAC
\end{minipage} \\
\midrule\noalign{}
\endhead
\bottomrule\noalign{}
\endlastfoot
Google & £34,204 & 75,830 & 4.04\% & £0.45 & 929 & 1.14\(\times\) &
£36.82 \\
TikTok & £12,577 & 24,561 & 1.81\% & £0.51 & 289 & 0.90\(\times\) &
£43.52 \\
Meta & £42,467 & 75,243 & 1.62\% & £0.56 & 931 & 0.82\(\times\) &
£45.61 \\
Other & £6,105 & 10,601 & 1.07\% & £0.58 & 117 & 0.62\(\times\) &
£52.18 \\
\end{longtable}

\subsubsection{Region, Product and Channel
Mix}\label{region-product-and-channel-mix}

North East = 20.1\% of revenue (Figure 8) --- the launch market.
Sourdough loaves dominate the product mix at 44\%; 69\% of orders are
mobile and organic search produces 32\% of revenue (Figure 9).

\begin{figure}[H]

{\centering \includegraphics[width=0.8\linewidth]{figures/fig4_region_mix} 

}

\caption{Regional revenue mix.}\label{fig:unnamed-chunk-8}
\end{figure}

\begin{figure}[H]

{\centering \includegraphics[width=0.85\linewidth]{figures/fig13_device_channel} 

}

\caption{Device and channel mix.}\label{fig:unnamed-chunk-9}
\end{figure}

\subsubsection{Seasonality, Subscriptions and A/B
Test}\label{seasonality-subscriptions-and-ab-test}

Christmas (Nov--Feb) is the revenue peak --- Feb alone delivers £18.7k
(Figure 10). Active subscriptions grew to \(\approx\) 560 by Feb 2026
with average monthly churn of 4.5\% (Figure 11). The hero-banner A/B
test (Figure 12) returns a +27.5\% relative uplift, z = 6.90, p
\textless{} 0.001 --- strongest defensible win in the dataset.

\begin{figure}[H]

{\centering \includegraphics[width=0.85\linewidth]{figures/fig11_seasonality} 

}

\caption{Month-of-year seasonality.}\label{fig:unnamed-chunk-10}
\end{figure}

\begin{figure}[H]

{\centering \includegraphics[width=0.85\linewidth]{figures/fig12_subscription_dynamics} 

}

\caption{Subscription dynamics.}\label{fig:unnamed-chunk-11}
\end{figure}

\begin{figure}[H]

{\centering \includegraphics[width=0.65\linewidth]{figures/fig5_ab_test_uplift} 

}

\caption{A/B test result with 95\% CIs.}\label{fig:unnamed-chunk-12}
\end{figure}

\subsubsection{Outliers and Productive
Opportunities}\label{outliers-and-productive-opportunities}

A scan of the dataset surfaces three statistical outliers that are also
commercial opportunities. Each is large enough to move MALC if acted on.

\begin{longtable}[]{@{}
  >{\raggedright\arraybackslash}p{(\linewidth - 8\tabcolsep) * \real{0.2000}}
  >{\raggedright\arraybackslash}p{(\linewidth - 8\tabcolsep) * \real{0.2000}}
  >{\raggedright\arraybackslash}p{(\linewidth - 8\tabcolsep) * \real{0.2000}}
  >{\raggedright\arraybackslash}p{(\linewidth - 8\tabcolsep) * \real{0.2000}}
  >{\raggedright\arraybackslash}p{(\linewidth - 8\tabcolsep) * \real{0.2000}}@{}}
\toprule\noalign{}
\begin{minipage}[b]{\linewidth}\raggedright
\#
\end{minipage} & \begin{minipage}[b]{\linewidth}\raggedright
Outlier observed
\end{minipage} & \begin{minipage}[b]{\linewidth}\raggedright
Magnitude
\end{minipage} & \begin{minipage}[b]{\linewidth}\raggedright
Why it matters
\end{minipage} & \begin{minipage}[b]{\linewidth}\raggedright
Productive action
\end{minipage} \\
\midrule\noalign{}
\endhead
\bottomrule\noalign{}
\endlastfoot
1 & \textbf{February 2026 revenue spike} --- single month at £18.7k
versus the trailing-12-month median of \(\approx\) £8.7k. & \(\approx\)
2.1\(\times\) median; +18\% on the next highest month. & Confirms that
the Christmas/winter window (Nov--Feb) is the structural peak. The brief
assumed Easter; the data rejects that prior. & Front-load \textbf{the
headline campaign 6--8 weeks before December} and shift paid spend share
into Nov--Jan; treat spring (Mar--Apr) as retention and test time. \\
2 & \textbf{Google CTR 4.04\% vs Meta 1.62\% and TikTok 1.81\%.} &
\(\approx\) 2.5\(\times\) the average. & Reveals \textbf{intent gap} ---
Google captures branded local search at far higher click efficiency than
discovery platforms. & Pivot the Q1 paid mix to \textbf{60\% Google /
25\% Meta / 15\% TikTok} and gate Meta/TikTok scale-up to ROAS clearing
1.0\(\times\). \\
3 & \textbf{A/B test uplift +27.5\% with p\textless0.001} --- more than
2.75\(\times\) the pre-registered MDE of 10\%. & z = 6.90, p \textless{}
0.001 on n = 96,700. & The treatment hero is one of the largest CVR wins
observed in a single PDP design test (Kohavi, Tang \& Xu, 2020). &
\textbf{Ship Variant B day-one} at launch; queue the next four
high-leverage A/B tests immediately (delivery threshold, subscription
discount, PDP layout, checkout form). \\
\end{longtable}

Two additional supporting outliers reinforce the picture:
\textbf{subscription share} has slipped from 26\% (mid-2024) to
\(\approx\) 14\% (early 2026) despite absolute growth --- a retention
concentration risk the AI recommender and Q4 loyalty tier are designed
to reverse. \textbf{Sourdough Loaf concentration} at 44\% of revenue is
a product-mix risk that the persona-led merchandising plan partly hedges
by elevating Celebration Cakes (21\%) and the Bread Box subscription
(18\%).

\subsection{Forecast, Budget Reallocation and
Roadmap}\label{forecast-budget-reallocation-and-roadmap}

Base scenario assumes \textbf{4\%/mo compound MALC growth} ---
triangulated from (i) trailing-12-month CAGR, (ii) UK SME DTC-food
benchmarks (Statista, 2025), (iii) bottom-up unit economics. Low/High
are \$\pm\$2pp.

\begin{figure}[H]

{\centering \includegraphics[width=0.85\linewidth]{figures/fig6_forecast_scenarios} 

}

\caption{12-month MALC forecast scenarios (baseline 430 at Feb 2026).}\label{fig:unnamed-chunk-13}
\end{figure}

\begin{longtable}[]{@{}
  >{\raggedright\arraybackslash}p{(\linewidth - 12\tabcolsep) * \real{0.1429}}
  >{\raggedright\arraybackslash}p{(\linewidth - 12\tabcolsep) * \real{0.1429}}
  >{\raggedright\arraybackslash}p{(\linewidth - 12\tabcolsep) * \real{0.1429}}
  >{\raggedright\arraybackslash}p{(\linewidth - 12\tabcolsep) * \real{0.1429}}
  >{\raggedright\arraybackslash}p{(\linewidth - 12\tabcolsep) * \real{0.1429}}
  >{\raggedright\arraybackslash}p{(\linewidth - 12\tabcolsep) * \real{0.1429}}
  >{\raggedright\arraybackslash}p{(\linewidth - 12\tabcolsep) * \real{0.1429}}@{}}
\toprule\noalign{}
\begin{minipage}[b]{\linewidth}\raggedright
Scenario
\end{minipage} & \begin{minipage}[b]{\linewidth}\raggedright
M+3
\end{minipage} & \begin{minipage}[b]{\linewidth}\raggedright
M+6
\end{minipage} & \begin{minipage}[b]{\linewidth}\raggedright
M+9
\end{minipage} & \begin{minipage}[b]{\linewidth}\raggedright
M+12
\end{minipage} & \begin{minipage}[b]{\linewidth}\raggedright
Revenue M+12
\end{minipage} & \begin{minipage}[b]{\linewidth}\raggedright
\(\Delta\) vs baseline
\end{minipage} \\
\midrule\noalign{}
\endhead
\bottomrule\noalign{}
\endlastfoot
Low (+2\%/mo) & 456 & 484 & 513 & 545 & £11,094 & +27\% \\
Base (+4\%/mo) & 484 & 544 & 612 & 688 & £14,008 & +60\% \\
High (+6\%/mo) & 512 & 610 & 726 & 865 & £17,615 & +101\% \\
\end{longtable}

\textbf{Budget reallocation rule} (Top-1 ROAS \(\times\) 1.30; Top-2
\(\times\) 1.05; Top-3 \(\times\) 0.85; bottom \(\times\) 0.55), applied
quarterly to the paid mix (Figure 14) --- boosts Google, cuts
sub-1.0\(\times\) ROAS channels.

\begin{figure}[H]

{\centering \includegraphics[width=0.85\linewidth]{figures/fig15_budget_reallocation} 

}

\caption{Recommended quarterly budget reallocation.}\label{fig:unnamed-chunk-14}
\end{figure}

\subsubsection{Retention and CLV
consideration}\label{retention-and-clv-consideration}

A simplified Pareto/NBD-lite CLV (Fader, Hardie \& Lee, 2005) yields a
12-month CLV of approximately \textbf{£55} per active customer at the
current 4.5\% monthly churn and a 55\% gross margin assumption. Reducing
churn by 1 percentage point (to 3.5\%) lifts CLV by approximately 29\%,
which is why the retention investments in Q3 and Q4 are core to the plan
rather than an afterthought.

\subsubsection{Quarterly roadmap}\label{quarterly-roadmap}

\begin{longtable}[]{@{}
  >{\raggedright\arraybackslash}p{(\linewidth - 6\tabcolsep) * \real{0.2500}}
  >{\raggedright\arraybackslash}p{(\linewidth - 6\tabcolsep) * \real{0.2500}}
  >{\raggedright\arraybackslash}p{(\linewidth - 6\tabcolsep) * \real{0.2500}}
  >{\raggedright\arraybackslash}p{(\linewidth - 6\tabcolsep) * \real{0.2500}}@{}}
\toprule\noalign{}
\begin{minipage}[b]{\linewidth}\raggedright
Quarter
\end{minipage} & \begin{minipage}[b]{\linewidth}\raggedright
Focus
\end{minipage} & \begin{minipage}[b]{\linewidth}\raggedright
MALC uplift
\end{minipage} & \begin{minipage}[b]{\linewidth}\raggedright
Budget
\end{minipage} \\
\midrule\noalign{}
\endhead
\bottomrule\noalign{}
\endlastfoot
Q1 & Shopify MVP launch, GA4 + Klaviyo wiring, email capture, A/B winner
& +5--8\% & £2,000 \\
Q2 & SEO/blog programme; Bread Box LP; subscription gifting test &
+8--12\% & £2,500 \\
Q3 & Budget reallocation by ROAS; PDP optimisation; retention flows &
+10--15\% & £3,000 \\
Q4 & Christmas campaign; loyalty tier; subscription gifting; review &
+15--20\% & £4,000 \\
\end{longtable}

\subsection{Measurement and
Governance}\label{measurement-and-governance}

A weekly cadence with quarterly reallocation suits a microenterprise
(Croll \& Yoskovitz, 2013). The governance flow in Figure 15 codifies
who owns what --- collect, model, visualise, decide --- with a feedback
loop from quarterly learnings back into event design and budget
weighting.

\begin{figure}[H]

{\centering \includegraphics[width=0.85\linewidth]{figures/fig14_governance} 

}

\caption{Measurement and decision governance flow.}\label{fig:unnamed-chunk-15}
\end{figure}

\textbf{Weekly dashboard specification.}

\begin{longtable}[]{@{}
  >{\raggedright\arraybackslash}p{(\linewidth - 8\tabcolsep) * \real{0.2000}}
  >{\raggedright\arraybackslash}p{(\linewidth - 8\tabcolsep) * \real{0.2000}}
  >{\raggedright\arraybackslash}p{(\linewidth - 8\tabcolsep) * \real{0.2000}}
  >{\raggedright\arraybackslash}p{(\linewidth - 8\tabcolsep) * \real{0.2000}}
  >{\raggedright\arraybackslash}p{(\linewidth - 8\tabcolsep) * \real{0.2000}}@{}}
\toprule\noalign{}
\begin{minipage}[b]{\linewidth}\raggedright
Tile
\end{minipage} & \begin{minipage}[b]{\linewidth}\raggedright
Metrics
\end{minipage} & \begin{minipage}[b]{\linewidth}\raggedright
Tools
\end{minipage} & \begin{minipage}[b]{\linewidth}\raggedright
Owner
\end{minipage} & \begin{minipage}[b]{\linewidth}\raggedright
Action threshold
\end{minipage} \\
\midrule\noalign{}
\endhead
\bottomrule\noalign{}
\endlastfoot
Acquisition & Sessions, share by channel & GA4 + Looker Studio & Founder
& Organic share \textless{} 30\% \(\Rightarrow\) SEO push \\
Conversion & CVR, AOV, cart abandonment & GA4 + Shopify & Founder & CVR
\textless{} 2.5\% \(\Rightarrow\) next A/B test \\
Paid efficiency & Spend, ROAS, CAC by platform & Google/Meta/TikTok APIs
& Agency / founder & ROAS \textless{} 1.0\(\times\) \(\Rightarrow\)
pause / cut \\
Retention & Active subs, churn, repeat rate & Klaviyo + Shopify &
Founder & Churn \textgreater{} 5\% \(\Rightarrow\) win-back flow \\
Content & Posts, engagements, clicks to site & Meta/TikTok APIs & Social
assistant & Engagement \textless{} 2\% \(\Rightarrow\) pillar review \\
\end{longtable}

\subsection{Leadership Dashboard --- Weekly KPI
View}\label{leadership-dashboard-weekly-kpi-view}

The assignment brief requires a \emph{measurement plan with events and a
weekly dashboard}. TAB's leadership view operationalises this through a
single, founder-readable dashboard built on the same metrics computed in
§3.4 and refreshed each Monday morning.

\textbf{Purpose.} The dashboard is intentionally narrow. It answers four
operational questions: \emph{(1)} is acquisition healthy (sessions,
channel share, CVR), \emph{(2)} is each pound of paid spend profitable
(ROAS, CAC vs AOV), \emph{(3)} is the customer base growing (MALC,
active subscriptions, churn), and \emph{(4)} are the seasonal and
experimentation levers performing on plan. Anything that does not answer
one of those four questions is deliberately excluded --- Few (2009)
notes that microenterprise dashboards fail when they accumulate
decorative metrics.

\textbf{Events captured.} Standard GA4 + Shopify events feed the
dashboard: \texttt{page\_view}, \texttt{view\_item},
\texttt{add\_to\_cart}, \texttt{begin\_checkout},
\texttt{add\_payment\_info}, \texttt{purchase}, \texttt{sign\_up}
(email), \texttt{subscribe\_start}, \texttt{subscribe\_cancel}. Paid
platforms contribute \texttt{spend}, \texttt{impressions},
\texttt{clicks}, \texttt{conversions} and \texttt{attributed\_revenue}
via their reporting APIs. Each event carries a UTM tag
(\texttt{utm\_source}, \texttt{utm\_medium}, \texttt{utm\_campaign}) so
revenue can be reconciled by channel.

\textbf{Cadence and ownership.} Daily auto-refreshed traffic and spend;
weekly review by the founder; monthly retention and cohort review;
quarterly scenario refresh and budget reallocation (Figure 14). Each
tile carries an explicit owner and an action threshold so the dashboard
is \emph{decision-shaped}, not just informational (Croll \& Yoskovitz,
2013).

\textbf{Format (Figure 16).} The companion
\texttt{TAB\ Dashboard\ (HTML)} is a single-file, offline-safe page
rendered entirely as inline SVG --- no external scripts or CDN
dependencies. The colour code follows the leadership reporting palette
used throughout this report: blue for MALC/revenue/primary, green for
on-target/positive (subscription starts, A/B winner), coral for
under-target/alert (sub-1.0\(\times\) ROAS, cancels, off-target KPIs),
and orange for secondary KPIs (Orders). Status pills next to each KPI
give a verdict in a single colour, satisfying the pre-attentive design
principle that hue should only appear where it carries meaning
(Cleveland \& McGill, 1984).

\begin{figure}[H]

{\centering \includegraphics[width=1\linewidth]{../DASHBOARD_COMPILED} 

}

\caption{Leadership dashboard snapshot. Top: callout and KPI strip with status pills. Middle: MALC trend, A/B test, paid ROAS, region mix. Bottom: seasonality, forecast, subscriptions, device split. Live version available in the project root as TAB Dashboard (HTML).}\label{fig:dashboard-snapshot}
\end{figure}

\subsection{Ethics, Privacy and
Accessibility}\label{ethics-privacy-and-accessibility}

\textbf{GDPR-by-default.} A consent banner (OneTrust or Shopify Customer
Privacy API) presents strictly-necessary, analytics and marketing cookie
categories as separate opt-ins (ICO, 2024). All email signups use
\textbf{double opt-in} through Klaviyo, with a one-click unsubscribe in
every send. A Data Subject Access Request (DSAR) workflow is documented
in Appendix B and handled within the 30-day statutory window. Personal
data is minimised (email, name, address, phone) and retained for 24
months post-purchase.

\textbf{Accessibility (WCAG 2.2 AA).} The site targets WCAG 2.2 AA (W3C,
2023). The brand palette meets a 4.5:1 contrast ratio; AI-generated alt
text is human-reviewed; keyboard navigation is tested at every release;
body type is 16px minimum with 1.5\(\times\) line-height. A monthly
Axe/WAVE audit and an annual screen-reader walkthrough with NVDA close
the loop.

\textbf{Responsible AI.} The AI layer is constrained to non-decisioning
tasks (content drafts, recommendations) with founder review before send
or publish. Customer-facing AI features (for example a future WhatsApp
assistant) are disclosed in plain English at the point of interaction,
per the EU AI Act transparency obligations (European Commission, 2024).

\newpage

\section{Recommendations and
Conclusion}\label{recommendations-and-conclusion}

\subsection{Five actionable
recommendations}\label{five-actionable-recommendations}

\begin{enumerate}
\def\labelenumi{\arabic{enumi}.}
\tightlist
\item
  \textbf{Adopt MALC as the North Star and instrument it in Week 1.}
  Every subsequent decision flows from this single number; without it
  the plan reverts to bets.
\item
  \textbf{Ship the AI-augmented Shopify MVP in eight weeks}, with the
  Variant B hero in production from day one (+27.5\% CVR uplift, p
  \textless{} 0.001).
\item
  \textbf{Reverse the conventional paid split: 60\% Google / 25\% Meta /
  15\% TikTok} until Meta and TikTok ROAS clears 1.0\(\times\).
  Reallocate quarterly using the ROAS-rank rule.
\item
  \textbf{Lead with owned channels} --- email, subscription, local SEO.
  CAC there is effectively zero versus £36--£52 on paid. Target 2,500
  email sign-ups by end of Q2.
\item
  \textbf{Plan around Christmas, not Easter.} The dataset rejects the
  brief's prior --- build the headline annual campaign for Nov--Feb and
  use spring for retention and A/B testing.
\end{enumerate}

\subsection{Next steps}\label{next-steps}

\begin{enumerate}
\def\labelenumi{(\roman{enumi})}
\tightlist
\item
  Approve the £11,500 12-month budget and the MVP scope; (ii) appoint a
  fractional digital partner to own the Shopify build and the weekly
  dashboard; (iii) freeze Variant B into the launch theme; (iv) start
  measuring MALC in Week 1; (v) reconvene at end-Q1 to refresh the
  forecast.
\end{enumerate}

\subsection{Conclusion}\label{conclusion}

The TAB dataset tells a clear, actionable story: a strong local brand
with a genuine product asymmetry, currently paying for paid traffic that
does not return its cost. A focused, mobile-first, AI-augmented Shopify
MVP --- with MALC at its centre and a measurement loop that reallocates
budget every quarter --- is the right next step. The numbers in this
report are sufficient to commit; the rest is execution.

\newpage

\section{References}\label{references}

Agresti, A. (2018) \emph{Statistical methods for the social sciences}.
5th edn. Boston: Pearson.

Baymard Institute (2024) \emph{E-commerce checkout usability report}.
Copenhagen: Baymard.

Cleveland, W.S. \& McGill, R. (1984) ``Graphical perception: theory,
experimentation, and application to the development of graphical
methods'', \emph{Journal of the American Statistical Association},
79(387), pp.~531--554.

Croll, A. \& Yoskovitz, B. (2013) \emph{Lean analytics: Use data to
build a better startup faster}. Sebastopol: O'Reilly.

European Commission (2024) \emph{Regulation (EU) 2024/1689 of the
European Parliament and of the Council on harmonised rules on artificial
intelligence (the AI Act)}. Brussels.

Fader, P.S., Hardie, B.G.S. \& Lee, K.L. (2005) ``\,`Counting your
customers' the easy way: an alternative to the Pareto/NBD model'',
\emph{Marketing Science}, 24(2), pp.~275--284.

Few, S. (2009) \emph{Now you see it: Simple visualization techniques for
quantitative analysis}. Burlingame: Analytics Press.

Hyndman, R.J. \& Athanasopoulos, G. (2021) \emph{Forecasting: Principles
and practice}. 3rd edn. Melbourne: OTexts.

Information Commissioner's Office (2024) \emph{Guidance on the use of
cookies and similar technologies}. Wilmslow: ICO.

Kannan, P.K. \& Li, H.A. (2017) ``Digital marketing: A framework, review
and research agenda'', \emph{International Journal of Research in
Marketing}, 34(1), pp.~22--45.

Kohavi, R., Tang, D. \& Xu, Y. (2020) \emph{Trustworthy online
controlled experiments: A practical guide to A/B testing}. Cambridge:
Cambridge University Press.

Recharge Payments (2024) \emph{State of subscription commerce report
2024}. Santa Monica: Recharge.

Reichheld, F.F. \& Schefter, P. (2000) ``E-loyalty: Your secret weapon
on the web'', \emph{Harvard Business Review}, 78(4), pp.~105--113.

Shopify Inc.~(2025) \emph{Shopify platform documentation}. Available at:
\url{https://help.shopify.com/} (accessed 12 May 2026).

Statista (2025) \emph{E-commerce in the United Kingdom --- statistics
and facts}. Hamburg: Statista.

Stripe (2025) \emph{Checkout optimization benchmarks: DTC food and
beverage}. San Francisco: Stripe.

Taylor, S.J. \& Letham, B. (2018) ``Forecasting at scale'', \emph{The
American Statistician}, 72(1), pp.~37--45.

W3C (2023) \emph{Web Content Accessibility Guidelines (WCAG) 2.2}.
Cambridge MA: W3C.

Wickham, H. et al.~(2019) ``Welcome to the Tidyverse'', \emph{Journal of
Open Source Software}, 4(43), 1686.

Wickham, H. \& Bryan, J. (2023) \emph{R for Data Science}. 2nd edn.
Sebastopol: O'Reilly.

Wirth, R. \& Hipp, J. (2000) ``CRISP-DM: Towards a standard process
model for data mining'', in \emph{Proceedings of the 4th International
Conference on the Practical Applications of Knowledge Discovery and Data
Mining}. Manchester, pp.~29--39.

\newpage

\section{Appendix A: Calculations and
Methodology}\label{appendix-a-calculations-and-methodology}

\subsection{Reproducibility}\label{reproducibility}

The analysis is reproducible end-to-end. From the project root in
RStudio:

\begin{Shaded}
\begin{Highlighting}[]
\FunctionTok{source}\NormalTok{(}\StringTok{"scripts/RUN\_ALL.R"}\NormalTok{)   }\CommentTok{\# runs scripts 01{-}04 then knits this Rmd}
\end{Highlighting}
\end{Shaded}

Each script writes a \texttt{.RData} snapshot to \texttt{scripts/} and a
\texttt{.png} file per figure to \texttt{figures/}. The full chain
(clean -\textgreater{} metrics -\textgreater{} visuals -\textgreater{}
forecast -\textgreater{} PDF) completes in approximately three minutes
on a standard laptop.

\subsection{Formula reference (R)}\label{formula-reference-r}

\begin{Shaded}
\begin{Highlighting}[]
\CommentTok{\# Funnel}
\NormalTok{cvr  }\OtherTok{\textless{}{-}}\NormalTok{ orders }\SpecialCharTok{/}\NormalTok{ website\_sessions          }\CommentTok{\# Conversion rate}
\NormalTok{aov  }\OtherTok{\textless{}{-}}\NormalTok{ revenue }\SpecialCharTok{/}\NormalTok{ orders                   }\CommentTok{\# Average order value}
\NormalTok{ctr  }\OtherTok{\textless{}{-}}\NormalTok{ clicks }\SpecialCharTok{/}\NormalTok{ impressions               }\CommentTok{\# Click{-}through rate}
\NormalTok{cpc  }\OtherTok{\textless{}{-}}\NormalTok{ spend }\SpecialCharTok{/}\NormalTok{ clicks                     }\CommentTok{\# Cost per click}
\NormalTok{roas }\OtherTok{\textless{}{-}}\NormalTok{ attributed\_revenue }\SpecialCharTok{/}\NormalTok{ spend         }\CommentTok{\# Return on ad spend}
\NormalTok{cac  }\OtherTok{\textless{}{-}}\NormalTok{ spend }\SpecialCharTok{/}\NormalTok{ conversions                }\CommentTok{\# Customer acquisition cost}

\CommentTok{\# North Star}
\NormalTok{malc }\OtherTok{\textless{}{-}}\NormalTok{ orders }\SpecialCharTok{\%\textgreater{}\%}
  \FunctionTok{group\_by}\NormalTok{(}\AttributeTok{month =} \FunctionTok{floor\_date}\NormalTok{(order\_date, }\StringTok{"month"}\NormalTok{)) }\SpecialCharTok{\%\textgreater{}\%}
  \FunctionTok{summarise}\NormalTok{(}\AttributeTok{malc =} \FunctionTok{n\_distinct}\NormalTok{(customer\_id))}

\CommentTok{\# Retention}
\NormalTok{churn\_rate }\OtherTok{\textless{}{-}}\NormalTok{ subscription\_cancels }\SpecialCharTok{/}\NormalTok{ active\_subscriptions}

\CommentTok{\# CLV (Pareto/NBD{-}lite)}
\NormalTok{clv }\OtherTok{\textless{}{-}}\NormalTok{ (aov }\SpecialCharTok{*}\NormalTok{ orders\_per\_cust\_month }\SpecialCharTok{*}\NormalTok{ gross\_margin) }\SpecialCharTok{/}\NormalTok{ churn\_rate}

\CommentTok{\# Forecast (geometric)}
\NormalTok{malc\_t1 }\OtherTok{\textless{}{-}}\NormalTok{ malc\_t }\SpecialCharTok{*}\NormalTok{ (}\DecValTok{1} \SpecialCharTok{+}\NormalTok{ monthly\_growth)}
\end{Highlighting}
\end{Shaded}

\subsection{A/B test inference}\label{ab-test-inference}

Under the null hypothesis \(H_0: p_B = p_A\), the two-proportion
z-statistic (Agresti, 2018) is:

\[z = \frac{p_B - p_A}{\sqrt{\hat{p}(1-\hat{p})\left(\tfrac{1}{n_A}+\tfrac{1}{n_B}\right)}}, \quad \hat{p} = \frac{x_A + x_B}{n_A + n_B}.\]

For the TAB hero-banner test: \(x_A = 1{,}485\), \(n_A = 51{,}884\),
\(x_B = 1{,}635\), \(n_B = 44{,}816\), giving \(z = 6.90\) and a
two-tailed \(p < 0.001\) --- significant at any conventional level.

\subsection{Limitations}\label{limitations}

\begin{enumerate}
\def\labelenumi{\arabic{enumi}.}
\tightlist
\item
  \textbf{Attribution.} ROAS uses the dataset's last-click
  \texttt{attributed\_revenue}. A multi-touch model would likely
  re-weight Meta upward and Google slightly down.
\item
  \textbf{MALC proxy.} Where \texttt{customer\_id} is missing, MALC is
  bounded below by unique-order counts; the true unique-customer count
  may be 5--10\% higher.
\item
  \textbf{Single A/B test.} The +27.5\% uplift is on one design pair;
  rolling tests should continue to validate.
\item
  \textbf{Seasonality.} Based on two annual cycles; a Bayesian 95\% band
  would be wider than the deterministic scenario range used here.
\end{enumerate}

\newpage

\section{Appendix B: Privacy and DSAR
Workflow}\label{appendix-b-privacy-and-dsar-workflow}

A Data Subject Access Request (DSAR) at TAB is handled as follows:

\begin{enumerate}
\def\labelenumi{\arabic{enumi}.}
\tightlist
\item
  The request arrives via
  \href{mailto:privacy@tab.example}{\nolinkurl{privacy@tab.example}} or
  the website contact form.
\item
  The founder verifies the requester's identity using order email plus
  the last four characters of post-code.
\item
  Shopify and Klaviyo data are exported within five working days.
\item
  The response is delivered as a password-protected PDF within the
  30-day statutory deadline.
\item
  The request is logged in a tracker for audit purposes.
\end{enumerate}

Deletion requests are processed within the same window and propagated to
all third-party processors (Klaviyo, Recharge, Stripe). The full
register is held in the secure operations Notion workspace.

\textbf{End of report.}

\end{document}
